\documentclass[12pt]{article}
%packages
\usepackage{amsmath, amsthm, amssymb}
\usepackage{graphicx}
\usepackage{enumerate}
\usepackage[utf8]{inputenc}
% \usepackage[spaced=0]{diffcoeff}
\usepackage[italic=true]{derivative}
\usepackage{tcolorbox}
\usepackage{physics2}
\usephysicsmodule{ab, braket}
\usepackage{siunitx}
\usepackage{fancyhdr}
\usepackage[margin=1in]{geometry}
\usepackage{tikz}
\usepackage{pgfplots} % for the axis environment
% \usepackage{fullpage}

%environments
\newtheorem*{question}{Q}
\newtheorem*{prop}{Proposition}

\newenvironment{quest}[1]
{\vspace{0.3cm}\noindent\textbf{Q(#1)}\em}{\vspace{0.3cm}}

\newenvironment*{Remark}{\noindent\textbf{Remarks}\begin{enumerate}}{\end{enumerate}\vspace{1em}\noindent}


\usepackage{hyperref}
\hypersetup{
    colorlinks=true,
    linkcolor=blue,
    filecolor=magenta,      
    urlcolor=cyan,
    pdftitle={Overleaf Example},
    pdfpagemode=FullScreen,
}


%commands
\newcommand{\e}{\varepsilon}
\newcommand{\R}{\mathbb{R}}
\newcommand{\Z}{\mathbb{Z}}
\newcommand{\N}{\mathbb{N}}
\newcommand{\C}{\mathbb{C}}
\newcommand{\Q}{\mathbb{Q}}


\newcommand{\grad}{\vec \nabla}
\newcommand{\LL}{\mathcal{L}}
\newcommand{\vphi}{\varphi}
\newcommand{\p}{\partial}



\newcommand{\eq}{\Leftrightarrow}
\newcommand{\clm}{\textbf{Claim }}
\newcommand{\rmks}{\noindent\textbf{Remarks}}

\newcommand{\levi}{\mathfrak{E}}


\makeatletter
     \newcommand\vb{\@ifstar\boldsymbol\mathbf}
     \newcommand\va[1]{\@ifstar{\vec{#1}}{\vec{\mathrm{#1}}}}
     \newcommand\vu[1]{%
       \@ifstar{\hat{{#1}}}{\hat{{#1}}}}
    % \renewcommand*{\vec}[][default]{def}
\makeatother


\newcommand{\ev}[1]{\langle{#1}\rangle}
\newcommand{\vecp}[1]{\vec{#1}\:'}
%\renewcommand\qedsymbol{QED}

%%%%%%%%%%%%%%%%%%%%%%
% Set up header/footer
\pagestyle{fancy}
\fancyhead[LO,L]{Vishwanath Wimalasena}
\setlength{\headheight}{14.5pt}
\fancyfoot[LO,L]{}
\fancyfoot[CO,C]{\thepage}
\fancyfoot[RO,R]{}
\renewcommand{\headrulewidth}{0.5pt}
\renewcommand{\footrulewidth}{0.4pt}
\renewcommand{\headrulewidth}{0.5pt}
%%%%%%%%%%%%%%%%%%%%%%



\tcbuselibrary{theorems}
\tcbuselibrary{breakable}
\tcbuselibrary{skins}

%environments
\newtheorem{theorem}{Theorem}[section]
\newtheorem{proposition}{Proposition}[section]
\newtheorem{definition}{Definition}[section]
\newtheorem{lemma}{Lemma}[section]
\newtheorem{lem}{Lemma}
\newtheorem*{ex}{Exercise}

\theoremstyle{definition}
\newtheorem*{sol}{Sol}

\newtcbtheorem[number within=section]{cthm}{Theorem}%
{colback=green!8!white, colframe=green!40!black,fonttitle=\bfseries, fontupper=\itshape}{mainthm}

\newtcbtheorem[number within=section]{cexercise}{Exercise}%
{colback=green!8!white, colframe=green!40!black,fonttitle=\bfseries, fontupper=\itshape}{mainthm}

\newtcbtheorem[number within=section]{cclaim}{Claim}%
{colback=green!8!white,colframe=green!40!black,fonttitle=\bfseries, fontupper=\itshape}{claim}

\tcolorboxenvironment{sol}{breakable, enhanced, left=5mm,
    before skip=10pt, after skip=10pt, frame  hidden,colback=green!6!white,
    borderline west={1pt}{0pt}{green!60!black}}

\tcolorboxenvironment{lem}{
    enhanced jigsaw,colback=green!8!white, colframe=green!40!black,breakable,before skip=10pt,after skip=10pt }

\tcolorboxenvironment{proof}{% `proof' from `amsthm' 
    blanker,breakable,left=5mm,
    before skip=10pt,after skip=10pt,
    borderline west={1pt}{0pt}{green!60!black}}

% \newtcolorbox{ansbox}{tcbox raise base, leftrule=0pt, arc=0pt, rightrule=0pt,toprule=0mm,bottomrule=0mm, colback=green!8!white, colframe=green!20!black}

% \newtcolorbox{background}[1][]{enhanced,
%     colback=green!5!white,colframe=green!40!black,
%     colbacktitle=green!40!black,fonttitle=\bfseries,
%     frame hidden,
%     title=Background,
%     boxrule=0pt,
%     titlerule style=green!70!black, #1}

\tcolorboxenvironment{ex}{enhanced jigsaw, colback=green!8!white, colframe=green!40!black,breakable,before skip=10pt,after skip=10pt}

\tcolorboxenvironment{question}{% `proof' from `amsthm' 
    blanker,colback=green,breakable,left=5mm,
    before skip=10pt,after skip=10pt,
    borderline west={1mm}{0pt}{green!40!black}}

% \tcbset{one-sided-color/.style={
%         breakable,
%         enhanced,
%         boxrule=0pt,
%         frame hidden,
%         borderline ={1pt}{0pt}{#1},
%         colback=#1!10,
%         coltitle=#1!200,
%         after title={\ },
%         attach title  to upper,
%     }
% }

\newtcolorbox{q}{one-sided-color=green!60!black, title=\textbf{Q}}

\newtcolorbox{outline}[1]{one-sided-color=blue!70!black, title=\textbf{#1}}

% \newtcolorbox{ans}{breakable,
%         enhanced,
%         boxrule=0pt,
%         frame hidden,
%         borderline west={1pt}{0pt}{green!70!black},
%         % borderline south={1pt}{0pt}{green!70!black},
%         % after title={\ },
%         % attach title  to upper,
%         top=1pt,
%         % boxsep=1pt,
%         colback=white,
%         coltitle=green!70!black!200}


% \newtcolorbox{q2}{breakable,
%         enhanced,
%         boxrule=0pt,
%         frame hidden,
%         borderline west={1pt}{0pt}{green!70!black},
%         colback=green!70!black!10, after title={Q}}

% \newtcolorbox{exc}[1]{one-sided-color=green!70!black, title=\textbf{(#1)}}
\definecolor{bordercolor}{rgb}{0.13,0.57,0.07}
\definecolor{boxcolor}{rgb}{0.28,0.69,0.28}

\newtcolorbox{exc}[1]{
        breakable,
        enhanced,
        frame hidden,
        borderline ={1pt}{0pt}{bordercolor},
        colback=boxcolor!10,
        coltitle=bordercolor,
        after title={\ },
        attach title  to upper,
        title=\textbf{Exercise #1},
        }

\newtcolorbox{ans}{
    breakable,
    enhanced,
    boxrule=0pt,
    frame hidden,
    borderline west={1pt}{0pt}{bordercolor},
    top=0pt,
    colback=white,
    coltitle=bordercolor,
    }
% \newcommand{\exc}[1]{\begin{exc}#1\end{exc}}

\fancyhead[CO,C]{\textbf{Basic Spin Dynamics}}
\fancyhead[RO,R]{Initial}

\newcommand{\HH}{\mathcal{H}}

\begin{document}
\begin{exc}{1a}
    Eigenstates of spin-1/2 in magnetic field. 
\end{exc}
\begin{ans}
    $\HH = -\gamma \vec B \cdot \vec S$. Orient axes so that $\vec B = B \hat z$. Then, $\HH = -\gamma BS_z \implies$ stationary states are eigenstates of $S_z$, i.e $\ket{\uparrow}, \ket{\downarrow}$ with energies $\HH \ket{\uparrow} = -\frac{\gamma \hbar B}{2}\ket{\uparrow}$ and $\HH \ket{\downarrow} = \frac{\gamma \hbar B}{2}\ket{\downarrow}$.
\end{ans}

% \begin{q2}
% \textbf{Q} What does `real space' mean for spin states?
% \end{q2}


\begin{exc}{1c}
    Spin 1/2 system in magnetic field $\vec B = (B_x, 0, B_z)$. 
\end{exc}
\begin{ans}
Expressing $\HH$ in matrix form w.r.t $S_z$ basis gives
    \[\HH = -\frac{\gamma \hbar}{2}\ab(B_x\begin{pmatrix}
        0 & 1 \\ 
        1 & 0 \\
    \end{pmatrix} + B_z\begin{pmatrix}
        1 & 0 \\ 
        0 & -1 \\
    \end{pmatrix}) = -\frac{\gamma \hbar}{2}\begin{pmatrix}
        B_z & B_x \\ 
        B_x & -B_z \\
    \end{pmatrix}\]
    Need to solve IVP 
    \[
    i\hbar \odv{\ket{\psi(t)}}{t} = \HH \ket{\psi(t)}, \qquad \ket{\psi(0)} = \ket{\uparrow}
    \]

    Express $\ket{\psi(t)} = \begin{pmatrix}
        a(t) \\ 
        b(t) \\
    \end{pmatrix}$ 
    so that we have the system of ODEs
    \begin{align*}
        \dot a &= \hat{\gamma}\ab(B_z a + B_x b), \qquad a(0) = 1 \\
        \dot b &= \hat{\gamma}\ab(B_x a - B_z b), \qquad b(0) = 0
    \end{align*}
    where $\hat \gamma = i\gamma / 2$. Differentiating again yields 
    \begin{align*}
        \ddot a &= \hat \gamma B_z \dot a + \hat \gamma B_x \dot b \\ 
        &= \hat \gamma B_z \ab(\hat \gamma B_z a + \hat \gamma B_x b) + \hat \gamma B_x \ab(\hat \gamma B_x a - \hat \gamma B_z b) \\
        &= \hat \gamma^2 |B|^2a \\
        &= -\frac{\gamma^2}{4}|B|^2 a
    \end{align*}
    where $|B|^2 = B_x^2 + B_z^2$. Similarly, $\ddot b = -\frac{\gamma^2}{4}|B|^2 b$
    These are SHM equations with solutions, 
    \begin{align*}
        a(t) &= C \cos\ab(\frac{\gamma|B|t}{2}) + D\sin\ab(\frac{\gamma|B|t}{2})\\ 
        b(t) &= E \cos\ab(\frac{\gamma|B|t}{2}) + F \sin\ab(\frac{\gamma|B|t}{2})
    \end{align*}
    Plugging in I.Cs gives 
    \begin{align*}
        a(0) &= C = 1\\ 
        b(0) &= E = 0\\
        \dot a(0) &= \hat \gamma B_z = D\frac{\gamma|B|}{2} \implies D = \frac{2\hat \gamma B_z}{\gamma |B|} = i\frac{B_z}{|B|} \\
        \dot b(0) &= \hat \gamma B_x = F\frac{\gamma|B|}{2} \implies F = \frac{2\hat \gamma B_x}{\gamma |B|} = i\frac{B_x}{|B|}
    \end{align*}
    So, the solution is 
    \begin{align*}
        \ket{\psi(t)} = \begin{pmatrix}
            \cos\ab(\frac{\gamma|B|t}{2}) + \frac{iB_z}{|B|}\sin\ab(\frac{\gamma|B|t}{2})\\
            \frac{iB_x}{|B|}\sin\ab(\frac{\gamma|B|t}{2})
        \end{pmatrix}
    \end{align*}
    Then, the probability of measuring spin up at time $t = \tau$ is 
    \begin{align*}
        |a(\tau)|^2 &= \ab|\cos\ab(\frac{\gamma|B|\tau}{2}) + \frac{iB_z}{|B|}\sin\ab(\frac{\gamma|B|\tau}{2})|^2\\ 
        &= \cos^2\ab(\frac{\gamma|B|\tau}{2}) + \frac{B_z^2}{|B|^2}\sin^2\ab(\frac{\gamma|B|\tau}{2})
    \end{align*}
    Note $B_x \to 0 \implies |B|^2 \to B_z^2 \implies |a(\tau)|^2 \to 1$. 
\end{ans}


% \section{Rotating Wave Approximation and Rabi oscillations}

\newpage 
\begin{exc}{2a}
    Write down the (time-dependent) Hamiltonian for the spin-1/2 system
\end{exc}
\begin{ans}
\[\HH = - \gamma \vb S \cdot \vb B = -\gamma[S_z B_0 + S_x B_1 \cos \nu t]\]
\end{ans}

\begin{exc}{2b}
    Derive ODE for system expressing $\ket \chi = a(t) \ket{0} + b(t) \ket{1}$ where $S_z \ket{0} = \hbar/2 \ket 0$ and $S_z \ket{1} = -\hbar/2 \ket 1$.
\end{exc}
\begin{ans}
    Have 
    \begin{align*}
        \HH \ket \chi 
        &= -\frac{\gamma \hbar}{2}(a(t)B_0 \ket 0 - b(t) B_0\ket 1 + a(t)B_1 \cos (\nu t)\ket 1  + b(t)B_1 \cos(\nu t)\ket 0)\\
        &= -\frac{\gamma \hbar}{2}((a(t)B_0 + b(t)B_1 \cos(\nu t)\ket 0 ) + (a(t)B_1 \cos(\nu t) - b(t)B_0)\ket 1 )
    \end{align*}
    Schrödinger gives $\HH \ket \chi = i\hbar \odv{\ket \chi}{t} = \dot a(t) \ket 0 + \dot b(t) \ket 1$. Setting equal to above and using orthonormality of basis, can decompose into 
    \begin{align}
     \dot a(t) &= i\frac{\omega}{2} a(t) + i\Omega \cos(\nu t)b(t)\\ 
    \dot b(t) &= -i \frac{\omega}{2} b(t) + i \Omega \cos(\nu t)a(t)
    \end{align}
    with 
    \[\boxed{\omega = \gamma B_0} \qquad \boxed{\Omega = \frac12\gamma B_1}\]
\end{ans}
\begin{exc}{2c}
    Derive using rotating frame $a(t) = \alpha(t)e^{i\nu t/2}, b(t) = \beta(t)e^{-i\nu t/2}$ and Rotating Wave Approximation $(\nu \gg \omega, \Omega)$.
\end{exc}
\begin{ans}
    \begin{align}
        \dot\alpha(t) &= i\frac{\delta}2\alpha(t) + i \frac{\Omega}2 \beta(t) \\ 
        \dot \beta (t) &= -i\frac{\delta}2\beta(t) + i \frac{\Omega}2 \alpha(t)
    \end{align}
\end{ans}
\begin{exc}{2d}
Solve system with IC $\ket{\chi(0)} = \ket 0$
\end{exc}
\begin{ans}
Equivalent to 1C, get 
\[\begin{pmatrix}
    \alpha(t) \\ 
    \beta(t) 
\end{pmatrix} = 
\begin{pmatrix}
    \displaystyle \cos\ab(\frac12 \Omega_R t) + i \frac{\delta }{\Omega_R} \sin\ab(\frac12 \Omega_R t) \\[1em] 
    \displaystyle i \frac{\Omega}{\Omega_R}\sin\ab(\frac12 \Omega_R t)
\end{pmatrix}
\]
\end{ans}
\begin{exc}{2e}
    Solve system at resonance $\delta = 0$ with $\Omega(t)$ time-dependent. 
\end{exc}
\begin{ans}
\begin{align*}
    \dot \alpha(t) &= \frac{i\Omega(t)}{2}\beta(t) \\ 
    \dot \beta(t) &= \frac{i\Omega(t)}{2}\alpha(t)
\end{align*}
Set $f = \alpha - \beta$ and $g = \alpha + \beta$. Then, we can express the system as 
\begin{align*}
    \dot f &= -\frac{i\Omega(t)}{2}f \\ 
    \dot g &= \frac{i\Omega(t)}{2}g
\end{align*}
Using $A(t) = \int^t \Omega(\tau)d\tau$, separating variables and solving gives 
\[f(t) = Ce^{-iA(t)/2}, \qquad g(t) = De^{iA(t)/2}\]
I.Cs give $f(0) = \alpha(0) - \beta(0) = 1$ and $g(0) = \alpha(0) + \beta(0) = 1$. So, $C = D = 1$. Returning to $\alpha, \beta$, we have 
\begin{align*}
    \alpha(t) &= \frac{1}{2}\ab(f(t) + g(t)) = \frac{1}{2}\ab(e^{-iA(t)/2} + e^{iA(t)/2}) = \cos\ab(\frac{A(t)}{2})\\
    \beta(t) &= \frac{1}{2}\ab(g(t) - f(t)) = \frac{1}{2}\ab(e^{iA(t)/2} - e^{-iA(t)/2}) = i\sin\ab(\frac{A(t)}{2})
\end{align*}
\end{ans}


\begin{exc}{2g}
    Write Hamiltonian for RWA ODE system 
\end{exc}
\begin{ans}
    \[i\hbar \odv{}{t}\begin{pmatrix}
        \alpha(t) \\ \beta(t)
    \end{pmatrix}
    = -\frac{\hbar}{2}\begin{pmatrix}
        \delta & \Omega \\ 
        \Omega & -\delta
    \end{pmatrix}
    \begin{pmatrix}
        \alpha(t) \\ \beta(t)
    \end{pmatrix}
    \]
\end{ans}

\begin{exc}{2h}
    Eigenstates of $\hat{\vb{H}}_{RWA}$? Eigenstates when $\Omega \gg \delta$ and $\delta \gg \Omega$ ? Plot eigenenergies for fixed $\Omega$. Energy splitting at $\delta = 0$?
\end{exc}
\begin{ans}
    Eigenstates found identically to question 1. Get 
    \[\ket{E_\pm} = \begin{pmatrix}
        \delta \pm \sqrt{\delta^2 + \Omega^2}\\
        \Omega
    \end{pmatrix}, \qquad E_{\pm} = \pm \frac12\sqrt{\delta^2 + \Omega^2}\]
    For $\delta \gg \Omega$, hamiltonian looks diagonal so expect basis states to be eigenstates. For $\Omega \gg \delta$, hamiltonian looks like $\sigma_x$ so expect symmetric and antisymmetric combinations of basis states to be eigenstates.

    Splitting at resonance is given by $E_+(\delta = 0) - E_-(\delta = 0) = \sqrt{\Omega^2} - (-\sqrt{\Omega^2}) = \Omega$
\end{ans}

\begin{exc}{2i}
     How does this entire analysis change if $\vec B_1||\vu y?$ Determine 
$\vb H_{RWA}$ for $\vec B_{AC} = B_1(\cos\theta, \sin\theta,0)\cos\nu t$. You
should find it behaves exactly the same as if $\vec B_{AC} = B_1(1, 0, 0)\cos(\vu t - \theta)$!
\end{exc}
\begin{proof}
    Recalling $S_y\ket{0} = i\frac\hbar 2\ket{1}$ and $S_y \ket{1} = -i\frac\hbar2\ket{0}$, we see that 
\begin{align*}
     \HH \ket \chi 
        &= -\frac{\gamma \hbar}{2}(a(t)B_0 \ket 0 - b(t) B_0\ket 1 + a(t)B_1 \cos (\nu t)\ket 1  + b(t)B_1 \cos(\nu t)\ket 0)\\
        &= -\frac{\gamma \hbar}{2}((a(t)B_0 + b(t)B_1 \cos(\nu t)\ket 0 ) + (a(t)B_1 \cos(\nu t) - b(t)B_0)\ket 1 )
\end{align*}

Setting up as before, we find 
\begin{align*}
    \dot a(t) &= \frac{i\omega}{2}a(t) + i\Omega\cos(\nu t)e^{-i\theta}b(t) \\ 
    \dot b(t) &= -\frac{i\omega}{2}b(t) + i\Omega\cos(\nu t)e^{-i\theta}a(t) 
\end{align*}
We can proceed as before in the rotating frame using $\alpha(t) = a(t)e^{i\nu t/2}$ and $\beta(t) = b(t)e^{-i\nu t/2}$ and get the equations 
\begin{align*}
    \dot \alpha(t) &= \frac{i\delta}{2}\alpha(t) + i\frac{\Omega}{2}(1 + e^{-2i\nu t})e^{-i\theta}\beta(t) \\ 
    \dot b(t) &= -\frac{i\delta}{2}\beta(t) + i\frac{\Omega}2 (1 + e^{2i\nu t})e^{i\theta}\alpha(t)
\end{align*}
Applying the rotatin wave approximation makes the $e^{\pm 2i \nu t} \to 0$, so our Hamiltonian looks like 
\begin{align*}
    \vb H_{RWA} = -\frac{\hbar}2\begin{pmatrix}
        \delta & \Omega e^{-i\theta} \\ 
        \Omega e^{i\theta} & -\delta
    \end{pmatrix}
\end{align*}

If instead we used $B_{AC} = B_1(1, 0, 0)\cos(\nu t - \theta)$, we would arrive at 

\begin{align*}
     \dot a(t) &= i\frac{\omega}{2} a(t) + i\Omega \cos(\nu t - \theta)b(t)\\ 
    \dot b(t) &= -i \frac{\omega}{2} b(t) + i \Omega \cos(\nu t - \theta)a(t)
\end{align*}

When we move to the rotating frame in the first equation, 
\[\cos(\nu t - \theta) = \frac12(e^{-i\theta} e^{i \nu t} + e^{i \theta}e^{-i \nu t}) \to \frac12(e^{-i\theta} e^{i \nu t} + e^{i \theta}e^{-i \nu t}) e^{-i\nu t} = \frac12(e^{-i\theta} + e^{i\theta}e^{-2i\nu t})\]
The RWA removes the latter term so that the same phase factor $e^{-i \theta}$ emerges. Similarly, we get the factor $e^{i\theta}$ in the bottom equation. 
\end{proof}

\begin{exc}{2j}
    What happens if $\vb B_1 \parallel \vu z$
\end{exc}
\begin{ans}
    THen our hamiltonian just looks like $\HH = -\gamma S_z (B_0 + B_1\cos\nu t)$ and we get zeeman splitting modulated by $B_1 \cos \nu t$.
    \[E_{\pm} = \mp \gamma \frac{\hbar}2 (B_0 + B_1\cos\nu t)\]
\end{ans}

\newpage 
\begin{exc}{3a}

\end{exc}

\begin{exc}{3b} 

\end{exc}
\begin{ans}
\begin{enumerate}[(i)]
    \item \[\hat U_{\pi_x/2} = \frac1{\sqrt 2}\begin{pmatrix}
        1 & i \\ i & 1
    \end{pmatrix}\]
    We see that $\hat U_{\pi/2}$ takes $\ket 0 \to \frac1{\sqrt 2}({\ket 0 + i\ket{1}}) = \ket{0_y}$ and $\ket 1 \to \frac1{\sqrt 2}(i \ket 0 + \ket 1) = \ket{1_y}$. It rotates the bloch vector by $\pi/2$ counter clockwise around the $x$-axis.

    \item \[\hat U_{\pi_x} = \frac1{\sqrt 2}\begin{pmatrix}
        0 & i \\ i & 0
    \end{pmatrix}\]
    Rotates around $x$ axis by $\pi$, takes $\ket 0 \leftrightarrow \ket 1$ 
\end{enumerate}
\end{ans}

\begin{exc}{3c}

\end{exc}
\begin{ans}
    \begin{enumerate}[(i)]
        \item \[\hat U_{\pi_y/2} = \frac1{\sqrt 2}\begin{pmatrix}
        1 & 1 \\ -1 & 1
    \end{pmatrix}\]
    Takes $\ket 0 \to \ket{1_x}$ and $\ket 1 \to \ket{0_x}$. Rotates by $\pi/2$ ccw around $y$ axis
    \item  \[\hat U_{\pi_y} = \frac1{\sqrt 2}\begin{pmatrix}
        0 & 1 \\ -1 & 0
    \end{pmatrix}\]
    $\pi$ turn around $y$ axis. Flips $\ket 0 \leftrightarrow \ket 1$
    \end{enumerate}
\end{ans}

\begin{exc}{3d}

\end{exc}
\begin{ans}
%     $\Omega \gg \delta \implies \Omega_R \approx \Omega$.
% \[\hat U_{HP} = \begin{pmatrix}
%     \cos(\Omega t/2) & i e^{-i \theta}\sin(\Omega t/2)\\ 
%     i e^{i \theta}\sin(\Omega t/2) & \cos(\Omega t /2)
% \end{pmatrix}\]
Then, 
Note 
\[\hat U(\Omega = 0,\delta,t) = \begin{pmatrix}
    e^{i\delta t/2} & 0 \\ 
    0 & e^{-i \delta t/2}
\end{pmatrix}\]
% \[\hat U_{HP, \pi_x/2} = \hat U(t, \delta = 0)\]
\begin{align*}
   \hat U_{Ramsey} &= \hat U_{\pi_x/2} \hat U(0, 0, t) \hat U_{\pi_x/2}\\ 
   &= \frac1{\sqrt 2}\begin{pmatrix}
        1 & i \\ i & 1
    \end{pmatrix}\begin{pmatrix}
       e^{i\delta t/2} & 0 \\
        0 & e^{-i\delta t/2}
    \end{pmatrix} \frac1{\sqrt 2}\begin{pmatrix}
        1 & i \\ i & 1
    \end{pmatrix}\\
    &= \frac12 \begin{pmatrix}
        1 & i \\ i & 1
    \end{pmatrix} \begin{pmatrix}
        e^{i\delta t/2} & i e^{i\delta t/2}\\
        i e^{-i\delta t/2} & e^{-i\delta t/2}
    \end{pmatrix}\\
    &= \begin{pmatrix}
       i \sin (\delta t/2) & i \cos(\delta t/2)\\ 
       i \cos(\delta t/2) & -i\sin(i \delta t/2) 
    \end{pmatrix}
\end{align*}
If state starts in $\ket 0 $, then after Ramsey Sequence $\ket 0 \to i \sin (\delta t/2)\ket 0 + i \cos(\delta t/2)\ket 1$ so there is probabliity $\cos^2(\delta t/2)$ that it is in $\ket 1$.

If we apply a $pi_y/2$ resonance pulse last instead, we get 
\begin{align*}
\hat U &= \hat U_{\pi_y/2} \hat U(0, 0, t) \hat U_{\pi_x/2}\\ 
   &= \frac1{\sqrt 2}\begin{pmatrix}
        1 & 1 \\ -1 & 1
    \end{pmatrix}\begin{pmatrix}
       e^{i\delta t/2} & 0 \\
        0 & e^{-i\delta t/2}
    \end{pmatrix} \frac1{\sqrt 2}\begin{pmatrix}
        1 & i \\ i & 1
    \end{pmatrix}\\
    &= \frac12 \begin{pmatrix}
        1 & 1 \\ -1 & 1
    \end{pmatrix} \begin{pmatrix}
        e^{i\delta t/2} & i e^{i\delta t/2}\\
        i e^{-i\delta t/2} & e^{-i\delta t/2}
    \end{pmatrix}
\end{align*}
The relevant matrix element for the transitiona amplitude $\ket 0 \to \ket 1$ is the lower left, which is given by 
\[\frac12(-e^{i\delta t/2} + ie^{-i\delta t/2})\]
The probability is then 
\begin{align*}
    P &= |\frac12(-\cos(\delta t/2) - i\sin(\delta t/2) + i\cos(\delta t/2) + \sin(\delta t/2))|^2\\
    &= \frac14((\sin \delta t/2 - \cos(\delta t/2))^2 + (\sin \delta t/2 - \cos \delta t/2)^2)\\
    &= \frac12(1 - 2\sin(\delta t/2)\cos(\delta t/2))\\
    &= \frac12(1 - \sin(\delta t))
\end{align*}
\end{ans}
\end{document}